\chapter{Sprint 5 --- Batch link + Knowledge/Support full BA (2026-05-16)}

\section{Bối cảnh}

Sprint 4.4 đã build skeleton cho \texttt{wujia\_portal\_knowledge} và \texttt{wujia\_portal\_support}, nhưng so với sheet BA \texttt{1. Model Field} (rows 117, 328--360, 372--430) còn thiếu nhiều field. Đồng thời \texttt{sale.order} chưa có field \texttt{batch\_id} để link đơn hàng với chuyến giao (BA row 117). Sprint 5 hoàn thiện 3 mảng song song:

\begin{itemize}[leftmargin=*]
    \item \textbf{Task A} --- thêm \texttt{sale.order.batch\_id} compute store index.
    \item \textbf{Task C} --- mở rộng knowledge full BA: state Selection, portal\_badge, tag M2m, sequence KNW-, cron expired.
    \item \textbf{Task E} --- mở rộng support ticket full BA: rename field, category M2o, 6-state workflow, analytics qua \texttt{message\_post}.
\end{itemize}

Decisions chốt với user: scope full BA (không cắt), batch\_id order theo id ASC, string mới English + \texttt{vi\_VN.po}, namespace giữ \texttt{wujia.*}, migration drop+init (skeleton chưa có data prod).

\section{Task A --- \texttt{sale.order.batch\_id}}

BA spec row 117: lấy picking của SO (\texttt{picking\_ids}) $\rightarrow$ batch\_id của picking đầu tiên không cancel.

\begin{lstlisting}[style=python]
batch_id = fields.Many2one(
    'stock.picking.batch',
    string='Delivery batch',
    compute='_compute_batch_id',
    store=True, index=True,
)

@api.depends('picking_ids.batch_id', 'picking_ids.state')
def _compute_batch_id(self):
    for order in self:
        pickings = order.picking_ids.filtered(
            lambda p: p.state != 'cancel' and p.batch_id
        ).sorted('id')
        order.batch_id = pickings[:1].batch_id
\end{lstlisting}

\textbf{Performance cho 1500 user portal}:

\begin{itemize}[leftmargin=*]
    \item Depend chỉ 2 field (\texttt{picking\_ids.batch\_id}, \texttt{picking\_ids.state}), không depend bare \texttt{picking\_ids} $\rightarrow$ giảm recompute thừa.
    \item \texttt{sorted('id')} in-memory trên O2m prefetched, không gọi \texttt{search()} $\rightarrow$ không thêm query khi compute trên list nhiều SO.
    \item \texttt{store=True + index=True} $\rightarrow$ filter "Đã có batch" và group\_by "Theo batch" trên list view nhanh ($O(\log n)$ qua btree).
    \item Không \texttt{compute\_sudo}: picking ACL đã được sale\_stock grant cho portal user qua record rule.
\end{itemize}

View update: form (sau \texttt{area\_id}), tree (optional hide), search (input + filter has\_batch + group\_by batch).

\section{Task C --- Knowledge full BA}

\subsection{Field changes \texttt{wujia.knowledge.article}}

\begin{longtable}{p{4cm}p{4cm}p{6cm}}
\toprule
Field & Action & Note \\
\midrule
\endhead
\texttt{code} & THÊM & Sequence \texttt{KNW-000001} qua \texttt{ir.sequence}. \\
\texttt{state} & THÊM (thay \texttt{published}) & Selection draft/published/archived, index + tracking. \\
\texttt{portal\_badge} & THÊM & Selection normal/new/important/mandatory. \\
\texttt{expired\_date} & THÊM & Datetime, index. \\
\texttt{sequence} & THÊM & Integer default 10, index (cho sort portal). \\
\texttt{is\_published\_portal} & THÊM compute store & Index --- portal domain filter 1 cột thay vì 2 điều kiện. \\
\texttt{attachment\_ids} & THÊM & M2m \texttt{ir.attachment}. \\
\texttt{active} & THÊM & Boolean default True (archive mềm). \\
\texttt{summary} & ĐỔI Char(250) $\rightarrow$ Text & Translate=True. \\
\texttt{tag\_ids} & ĐỔI Char CSV $\rightarrow$ M2m & Model mới \texttt{wujia.knowledge.tag}. \\
\texttt{published} & XOÁ & Replaced by \texttt{state}. \\
\bottomrule
\end{longtable}

Inherit \texttt{mail.thread} + \texttt{mail.activity.mixin} cho tracking state.

\subsection{Model mới \texttt{wujia.knowledge.tag}}

Fields: \texttt{name} (unique, translate), \texttt{color} (Integer), \texttt{active}, \texttt{article\_count} (compute qua \texttt{\_read\_group} batch không vòng for query).

\subsection{Cron \texttt{\_cron\_recompute\_is\_published}}

\texttt{is\_published\_portal} depend trên \texttt{expired\_date > now} không tự re-trigger theo thời gian. Cron daily 03:00 re-evaluate cho article vừa hết hạn:

\begin{lstlisting}[style=python]
@api.model
def _cron_recompute_is_published(self):
    now = fields.Datetime.now()
    expired = self.search([
        ('is_published_portal', '=', True),
        ('expired_date', '!=', False),
        ('expired_date', '<=', now),
    ])
    if expired:
        expired._compute_is_published_portal()
\end{lstlisting}

\section{Task E --- Support Ticket full BA}

\subsection{ADR-016 --- Không tạo \texttt{wujia.support.ticket.message} riêng}

BA gợi ý model song song \texttt{franchise.support.ticket.message} cho author\_role snapshot / is\_internal flag / message\_type. Quyết định KHÔNG implement vì:

\begin{itemize}[leftmargin=*]
    \item \texttt{mail.message} của Odoo đã có \texttt{author\_id}, \texttt{body}, \texttt{attachment\_ids}, \texttt{message\_type}, \texttt{subtype\_id} (cho is\_internal qua \texttt{mail.mt\_note}).
    \item Tạo model song song duplicate hệ thống chatter $\rightarrow$ dual write phức tạp, dễ inconsistent.
    \item Analytics BA cần (\texttt{last\_message\_date}, \texttt{last\_response\_by}, \texttt{need\_customer\_reply}, \texttt{message\_count}) đặt ở level ticket --- update atomic qua \texttt{message\_post()} override.
\end{itemize}

YAGNI: nếu sprint sau BA yêu cầu report message-level với role snapshot bất biến, lúc đó mới tạo model riêng.

\subsection{Field changes \texttt{wujia.support.ticket}}

\textbf{Rename}: \texttt{subject} $\rightarrow$ \texttt{title}; \texttt{user\_id} $\rightarrow$ \texttt{created\_by\_id}; \texttt{handler\_user\_id} $\rightarrow$ \texttt{assigned\_user\_id}.

\textbf{Đổi kiểu}: \texttt{description} Text $\rightarrow$ Html(sanitize); \texttt{category} Selection $\rightarrow$ \texttt{category\_id} Many2one tới model mới \texttt{wujia.support.category}; \texttt{priority} 4-level $\rightarrow$ 2-level (normal/urgent).

\textbf{State 4 $\rightarrow$ 6 giá trị}: new / in\_progress / \textbf{waiting\_customer} / resolved / closed / \textbf{cancelled}.

\textbf{Thêm}: \texttt{created\_by\_role} (snapshot owner/manager/staff), \texttt{store\_partner\_id} (related store), \texttt{sale\_order\_id}, \texttt{picking\_batch\_id}, \texttt{assigned\_team\_id} (crm.team), \texttt{last\_message\_date}, \texttt{last\_response\_by}, \texttt{need\_customer\_reply}, \texttt{message\_count}, \texttt{attachment\_count}, \texttt{open\_date / in\_progress\_date / resolved\_date / closed\_date}, \texttt{cancel\_reason}, \texttt{internal\_note}, \texttt{portal\_visible}, \texttt{active}.

\subsection{\texttt{message\_post} override --- analytics atomic}

\begin{lstlisting}[style=python]
def message_post(self, **kwargs):
    msg = super().message_post(**kwargs)
    if msg and msg.message_type == 'comment':
        author = msg.author_id
        is_portal = bool(author.user_ids and author.user_ids[:1].share)
        vals = {'last_message_date': msg.date or fields.Datetime.now()}
        if is_portal:
            vals.update({
                'last_response_by': 'customer',
                'need_customer_reply': False,
            })
        else:
            vals['last_response_by'] = 'hq'
            vals['need_customer_reply'] = self.state == 'waiting_customer'
        self.sudo().write(vals)
    return msg
\end{lstlisting}

\subsection{Constraint cancel\_reason}

\begin{lstlisting}[style=python]
@api.constrains('state', 'cancel_reason')
def _check_cancel_reason(self):
    for rec in self:
        if rec.state == 'cancelled' and not (rec.cancel_reason or '').strip():
            raise ValidationError(_(
                "Ticket %s requires a cancel reason when state is 'cancelled'.",
                rec.name or rec.title or '?',
            ))
\end{lstlisting}

\subsection{Model mới \texttt{wujia.support.category}}

Tách category Selection thành master data: \texttt{name} unique, \texttt{code} unique, \texttt{default\_assigned\_user\_id}, \texttt{default\_assigned\_team\_id}, \texttt{sequence}, \texttt{active}. 7 default category seed qua \texttt{data/wujia\_support\_category\_data.xml} (order/delivery/product/pos/operation/account/other) --- khớp code với Selection cũ để controller portal mapping được.

\section{Migration --- bài học data loss}

Khi rename/đổi kiểu field trên model có data, Odoo 19 \texttt{-u} sẽ DROP cột cũ aggressively (mất data trong cột đó). Lần đầu chạy \texttt{-u} trên dev DB \texttt{wujia\_tea\_19} bị mất data ticket + article hiện có (subject/user\_id/handler\_user\_id/category/published bị drop).

Bài học: trước khi commit code có rename field trên model có data, hoặc

\begin{itemize}[leftmargin=*]
    \item Viết \texttt{<module>/migrations/<new\_ver>/pre-migrate.py} copy data sang cột mới trước khi Odoo drop.
    \item Hoặc xác nhận với user drop+init + reseed (chỉ khi data thuần test/skeleton).
\end{itemize}

Sprint này chốt phương án drop+init vì cả dev và prod đều là skeleton. Quy trình:

\begin{lstlisting}[style=shell]
# Trên dev (đã chạy)
bash /home/huyban/odoo-dev/WujiaTea/scripts/reseed_full.sh

# Trên prod server (sau khi git pull code mới)
# Xem chi tiết: docs/DEPLOY_SPRINT5.md
\end{lstlisting}

\section{Files thay đổi}

\begin{lstlisting}[style=shell]
SỬA:
  custom/wujia_sale/__manifest__.py                              # bump 19.0.3.0.0
  custom/wujia_sale/models/sale_order.py                         # +batch_id
  custom/wujia_sale/views/sale_order_views.xml                   # batch_id 3 views

  custom/wujia_portal_knowledge/__manifest__.py                  # +mail depend
  custom/wujia_portal_knowledge/models/__init__.py               # +tag
  custom/wujia_portal_knowledge/models/wujia_knowledge_article.py # full BA spec
  custom/wujia_portal_knowledge/models/wujia_knowledge_category.py # +description
  custom/wujia_portal_knowledge/security/ir.model.access.csv     # +tag rows
  custom/wujia_portal_knowledge/views/wujia_knowledge_backend_views.xml
  custom/wujia_portal_knowledge/views/portal_knowledge.xml       # render badge
  custom/wujia_portal_knowledge/controllers/portal.py            # is_published_portal filter

  custom/wujia_portal_support/__manifest__.py                    # +sale,batch,team
  custom/wujia_portal_support/models/__init__.py                 # +category
  custom/wujia_portal_support/models/wujia_support_ticket.py     # full BA spec rewrite
  custom/wujia_portal_support/security/ir.model.access.csv       # +category rows
  custom/wujia_portal_support/security/wujia_support_rules.xml   # rule split portal/HQ
  custom/wujia_portal_support/views/wujia_support_backend_views.xml # rewrite full
  custom/wujia_portal_support/views/portal_support.xml           # rename refs + chatter
  custom/wujia_portal_support/controllers/portal.py              # dynamic category + reply

  scripts/seed_portal_demo.py                                    # align schema mới

TẠO:
  custom/wujia_portal_knowledge/models/wujia_knowledge_tag.py
  custom/wujia_portal_knowledge/data/ir_sequence_data.xml         # seq KNW-
  custom/wujia_portal_knowledge/data/ir_cron_data.xml             # cron daily

  custom/wujia_portal_support/models/wujia_support_category.py
  custom/wujia_portal_support/data/wujia_support_category_data.xml # 7 cat

  custom/wujia_sale/i18n/{wujia_sale.pot, vi_VN.po}
  custom/wujia_portal_knowledge/i18n/{*.pot, vi_VN.po}
  custom/wujia_portal_support/i18n/{*.pot, vi_VN.po}

  scripts/seed_knowledge_demo.py
  scripts/seed_support_demo.py
  scripts/test_sprint5.py
  scripts/reseed_full.sh

  docs/DEPLOY_SPRINT5.md
\end{lstlisting}

\section{Verify (đã pass)}

\texttt{scripts/test\_sprint5.py} chạy 20 assertion: \textbf{20 PASS / 0 FAIL} (1 SKIP cho batch\_id vì legacy seed\_portal\_demo chưa tạo SO --- liên quan thiếu product seed, không phải bug Sprint 5).

Check covered:

\begin{itemize}[leftmargin=*]
    \item Knowledge: published article có \texttt{is\_published\_portal=True}; expired article bị exclude; tag M2m search OK; sequence KNW- assigned; cron exists.
    \item Support: 7 default category seeded; ticket có title/category\_id/created\_by\_id; priority đúng BA spec; cancel/waiting\_customer states tồn tại; cancel\_reason required (ValidationError); state transition stamp dates; \texttt{message\_post} cập nhật \texttt{last\_message\_date} + \texttt{last\_response\_by}.
\end{itemize}

% ===========================================================================
